% GENERATED FILE. Edit canonical lessons, then run npm run generate:books.
% canonical-source-hash: fnv1a64:44cca1a9c8438a54
% canonical-lessons: GU-C31-nathi, GU-C31-samajto-nathi, GU-W08-pha, GU-C31-maaf, GU-C31-repair-five, GU-R35-second-pass-read-it-or-repair-it

\chapter{The Word for Not, and the Word for Sorry}
\label{ch:gu-not-and-sorry}
\hlchaptermodality{35}

\begin{chapteropening}
\textbf{By the end of this chapter:} \emph{I can say that something is not so, tell someone I did not understand, and apologise or get a stranger's attention.}

The repair kit. Gujarati negates a sentence by swapping chhe for nathi rather than adding a word, and that one swap unlocks the sentence a beginner needs most -- samajto nathi, 'I do not understand'. The chapter spends the only new letter in this whole run, pha, and spends it on maaf karo, which both apologises and stops a stranger to ask.
\end{chapteropening}

\section[nathī]{\gu{નથી} (nathī) — is not}

\label{lesson:GU-C31-nathi}



\begin{quote}
Two things you have had since the first chapters. Say the Gujarati for \textbf{yes} and for \textbf{no}. Then say the little word that means \textbf{is}.

\begin{itemize}
\raggedright
  \item \emph{Recall:} \emph{hā}, \emph{nā}, and \emph{chhe}
\end{itemize}

You can agree and you can refuse. What you cannot yet do is say that something \textbf{is not the case} — and until you can, you cannot say you did not understand.
\end{quote}

\begin{cousinweb}
Gujarati does not negate \textbf{\gu{છે}} by putting a word in front of it. It swaps the whole word out:

\begin{quote}
\textbf{\gu{છે}} — is · \textbf{\gu{નથી}} — is not
\end{quote}

That is the entire rule, and it is a rule about one word rather than about sentences: wherever \emph{chhe} stands, \emph{nathī} can stand instead.

The two halves are still visible inside it. \textbf{\gu{ન}} is the old \emph{na}, \textquotedblleft{}not\textquotedblright{}. The \textbf{-\gu{થી}} is what is left of \textbf{\gu{અસ્તિ}} (\emph{asti}), \textquotedblleft{}is\textquotedblright{} — the same ancient verb that comes down into English as \textbf{is} and into Latin as \emph{est}. So \emph{nathī} is literally \textbf{not-is}, two words that fused so long ago that Gujarati now treats them as one.
\end{cousinweb}

\subsection*{Your turn}
\begin{itemize}
\raggedright
  \item \emph{Say it:} \emph{chhe} — then \emph{nathī}. One swap, nothing added.
  \item \emph{Contrast:} \emph{chhe} / \emph{nathī}, twice, until the pair is automatic.
  \item \emph{Say it:} \emph{nathī} alone, as a flat answer. It is a whole reply.
\end{itemize}

\begin{tcolorbox}[breakable,colback=teal!4,colframe=teal!35!black,title={Before you move on}]
How does Gujarati say \textbf{is not}? (\textbf{\gu{નથી}} — it replaces \emph{chhe} rather than joining it.) What two old words are hiding inside it? (\emph{na}, \textquotedblleft{}not\textquotedblright{}, and \emph{asti}, \textquotedblleft{}is\textquotedblright{}.) Next: the sentence a learner needs more than any other.
\end{tcolorbox}
\section[samajto nathī]{\gu{સમજતો} \gu{નથી} (samajto nathī) — I do not understand}

\label{lesson:GU-C31-samajto-nathi}



\begin{quote}
One lesson back: the word that replaces \emph{chhe}. And from the verb chapter, long before this one: the verb \textbf{to understand}.

\begin{itemize}
\raggedright
  \item \emph{Recall:} \emph{nathī}, and \emph{samajvũ}
\end{itemize}
\end{quote}

\begin{cousinweb}
Put the verb into the form that describes what you are doing, and follow it with the negator:

\begin{quote}
\textbf{\gu{સમજતો} \gu{નથી}} — \emph{samajto nathī} — \textquotedblleft{}(I) am not understanding\textquotedblright{}
\end{quote}

Gujarati leaves out the \textbf{\gu{હું}} here, the way English leaves out the subject of \emph{don't know}. The sentence is complete without it.

\textbf{One thing to know before you say it.} The \textbf{-\gu{તો}} ending is the masculine form. A woman says \textbf{\gu{સમજતી} \gu{નથી}} (\emph{samajtī nathī}), with \textbf{-\gu{તી}}. Gujarati marks the speaker's own gender on this verb, so the sentence you use is the one that matches you — the same choice the track's gender chapter set up, arriving now on a sentence you will say every day.
\end{cousinweb}

\begin{culture}
This is the repair sentence. When a conversation runs ahead of you, this is what buys you the next attempt — and a beginner who cannot say it has to leave the conversation instead. It is worth more than any noun in this book.
\end{culture}

\subsection*{Your turn}
\begin{itemize}
\raggedright
  \item \emph{Say it:} \emph{samajto nathī} — or \emph{samajtī nathī}, whichever is yours.
  \item \emph{Say it:} it slowly, then at speed. Speed is the point; it has to arrive fast.
  \item \emph{Contrast:} \emph{nathī} alone, then the whole sentence.
\end{itemize}

\begin{tcolorbox}[breakable,colback=teal!4,colframe=teal!35!black,title={Before you move on}]
Say that you do not understand, in the form that matches you. Which letter changes between the two forms, and what does it mark? (\textbf{-\gu{તો}} against \textbf{-\gu{તી}} — the speaker's own gender.)
\end{tcolorbox}
\section[pha]{\gu{ફ} — the consonant pha}

\label{lesson:GU-W08-pha}



\begin{quote}
\begin{itemize}
\raggedright
  \item \emph{Read it:} \textbf{\gu{પ}}, then \textbf{\gu{મ}}
\end{itemize}

Both are already yours. Today adds one shape, and it is the only new letter in this whole run of chapters — everything else ahead is spelled with signs you have had for a long time.
\end{quote}

\subsection*{Script — meet one form}
Keep \textbf{\gu{ફ}} visible beside \textbf{\gu{પ}}. They are relatives and they are not the same: \emph{pa} is the plain sound, \emph{pha} is the same sound with a push of breath behind it. The shape says so too — \textbf{\gu{ફ}} carries an extra stroke that \textbf{\gu{પ}} does not.

Set them side by side once and look at the difference rather than at the whole:

\begin{quote}
\textbf{\gu{પ}} — pa · \textbf{\gu{ફ}} — pha
\end{quote}

\subsection*{Writing — notice and trace}
\hlblockfigure{\detokenize{figures/GU-W08-pha-filmstrip.pdf}}{How \gu{ફ} is written, stroke by stroke}

Finger-trace \textbf{\gu{ફ}} once, then trace it once in the air. Say \emph{pha} with the breath audible as your finger finishes.

\subsection*{Writing — guided copy, model in view}
Keep the model in view and copy \textbf{\gu{ફ}} once. Compare it against \textbf{\gu{પ}} beside it and repair only the stroke that separates them. Do not start a row.

\begin{tcolorbox}[breakable,colback=teal!4,colframe=teal!35!black,title={Before you move on}]
Which of \textbf{\gu{પ}} and \textbf{\gu{ફ}} carries the extra breath, and how does the written shape tell you? (\textbf{\gu{ફ}} — it carries the extra stroke.) The next lesson is the word this letter was bought for.

Source: \href{https://www.unicode.org/charts/PDF/U0A80.pdf}{Unicode Gujarati chart}.
\end{tcolorbox}
\section[māf karo]{\gu{માફ} \gu{કરો} (māf karo) — excuse me}

\label{lesson:GU-C31-maaf}



\begin{quote}
One lesson back, the new letter. And from the very first chapter, the word for \textbf{thank you}.

\begin{itemize}
\raggedright
  \item \emph{Recall:} \textbf{\gu{ફ}}, and \emph{ābhār}
\end{itemize}

You have been able to thank someone since page one of this book. You have never been able to apologise to them.
\end{quote}

\begin{cousinweb}
\begin{quote}
\textbf{\gu{માફ} \gu{કરો}} — \emph{māf karo} — \textquotedblleft{}do forgiveness\textquotedblright{}, that is, \textbf{forgive me}
\end{quote}

\textbf{\gu{માફ}} is the forgiveness itself, and \textbf{\gu{કરો}} is \emph{do} — the polite form you already use for asking. Gujarati does not have a verb \textquotedblleft{}to apologise\textquotedblright{}; it asks the other person to \textbf{do} the forgiving.

Read the first word slowly. \textbf{\gu{મ}} and the \textbf{\gu{ા}} stroke you have had for a long time; \textbf{\gu{ફ}} you learned in the last lesson. Three signs, and only one of them was new.

\textbf{\gu{માફ}} is a traveller. It came from Arabic \emph{muāf} by way of Persian, on the same road that carried \textbf{\gu{બજાર}} into the market chapter and \textbf{\gu{કાગળ}} onto the page two chapters ago. The apology in this book arrived the way the paper did.
\end{cousinweb}

\begin{culture}
\emph{Māf karo} does two jobs, and the second is the one a traveller needs most. It apologises — and it is also how you \textbf{get someone's attention} before you ask them anything: said to a stranger on the street, it means \emph{excuse me}, not \emph{I have wronged you}. One phrase opens the conversation and repairs it.
\end{culture}

\subsection*{Your turn}
\begin{itemize}
\raggedright
  \item \emph{Say it:} \emph{māf karo} — as an apology.
  \item \emph{Say it:} \emph{māf karo} again — this time to stop a stranger and ask something.
  \item \emph{Write it:} \textbf{\gu{માફ}}, with the model in view. One copy.
\end{itemize}

\begin{tcolorbox}[breakable,colback=teal!4,colframe=teal!35!black,title={Before you move on}]
Say \emph{excuse me} to a stranger. What does \textbf{\gu{કરો}} contribute, literally? (\emph{Do} — you are asking them to \textbf{do} the forgiving.) Which sign in \textbf{\gu{માફ}} did you learn one lesson ago? (\textbf{\gu{ફ}}.)
\end{tcolorbox}
\section[Practice]{The repair kit, all together}

\label{lesson:GU-C31-repair-five}



\begin{quote}
No new words. One run through the four things this chapter added, scored as a single performance.
\end{quote}

\begin{grammarlens}[title={What you've built}]
\begin{itemize}
\raggedright
  \item \textbf{\gu{નથી}} — \emph{is not}, which replaces \emph{chhe} rather than joining it.
  \item \textbf{\gu{સમજતો} / \gu{સમજતી} \gu{નથી}} — \emph{I do not understand}, in the form that matches you.
  \item \textbf{\gu{ફ}} — one new letter, the only one this run of chapters spends.
  \item \textbf{\gu{માફ} \gu{કરો}} — \emph{excuse me}, which both apologises and opens a conversation.
\end{itemize}

Before this chapter a reader of this book could greet, count, name a hundred things and ask for tea — and had no way at all to say that something was \textbf{not} so, or that they had not followed. The four items above are what a learner reaches for when everything else has failed.
\end{grammarlens}

\subsection*{Your turn}
\begin{itemize}
\raggedright
  \item \emph{Say it:} someone asks you something too fast. Two sentences: stop them, then tell them you did not follow.
  \item \emph{Say it:} \emph{nathī} alone, as a flat contradiction.
  \item \emph{Read it:} \textbf{\gu{માફ}} — name its three signs in order, then read the word.
  \item \emph{Say it:} the whole kit in order — \emph{māf karo}, \emph{samajto nathī}, \emph{nathī}.
\end{itemize}

\begin{tcolorbox}[breakable,colback=teal!4,colframe=teal!35!black,title={Before you move on}]
You are lost in a conversation. Say the two sentences that get you out, in the order you would really use them. (\emph{Māf karo} to stop them, then \emph{samajto nathī}.) Next chapter: the word for \textbf{and}, which this book has also never had.
\end{tcolorbox}
\section[māf karo]{Read it, or repair it}

\label{lesson:GU-R35-second-pass-read-it-or-repair-it}



\begin{quote}
You did not follow what was said. Two words, from memory, that buy you a second attempt.
\end{quote}

\subsection*{You'll want to know: four words you can now read unaided}
\noindent
\begin{tabularx}{\linewidth}{@{}>{\raggedright\arraybackslash}X>{\raggedright\arraybackslash}X>{\raggedright\arraybackslash}X@{}}
\toprule
\textbf{word} & \textbf{} & \textbf{} \\
\midrule
\textbf{\gu{નમસ્તે}} & \emph{namaste} & the greeting \\
\textbf{\gu{કાગળ}} & \emph{kāgaḷ} & paper \\
\textbf{\gu{માફ} \gu{કરો}} & \emph{māf karo} & forgive me \\
\textbf{\gu{નાક}} & \emph{nāk} & nose \\
\bottomrule
\end{tabularx}

Not one of them needed a new letter when it was read. Every piece was already yours, taught pages or chapters earlier and waiting — which is the whole reason the script came in a drizzle rather than a block.

\subsection*{Script: the letters that stand alone, and the two with breath}
\begin{quote}
\gu{એ}  \gu{ઉ}  \gu{ઈ}   and   \gu{ખ}  \gu{ઢ}
\end{quote}

Two groups, and the difference is worth naming. \textbf{\gu{એ}}, \textbf{\gu{ઉ}} and \textbf{\gu{ઈ}} are vowels standing on their own, needed only where no consonant is carrying them. \textbf{\gu{ખ}} and \textbf{\gu{ઢ}} are consonants with the breath let out after them — the same step that separates a plain letter from its aspirated partner all the way down the row.

Standing alone is a question of \textbf{position}. Aspiration is a question of \textbf{sound}. Neither tells you anything about the other.

\begin{grammarlens}[title={What you've built — courtesy, farewells, and the kit}]
\textbf{\gu{સારું}} — good, okay, fine — is the small word that keeps a conversation moving when you have nothing else to add. Beside it sit the courtesy exchange and the farewells, each practised once in its own chapter and not since.

And behind all of them, the repair kit this chapter finished: the way to say you did not understand, and the way to ask forgiveness for it. \textbf{A beginner uses the repair kit more than any single word in the book}, which is the argument for running it cold rather than only where it was taught.
\end{grammarlens}

\begin{tcolorbox}[breakable,colback=teal!4,colframe=teal!35!black,title={Before you move on}]
When does a Gujarati vowel need a letter of its own? (\textbf{When no consonant is carrying it.}) What separates \textbf{\gu{ખ}} from its plain partner? (\textbf{The breath after it.}) Which small word keeps a conversation moving? (\textbf{\gu{સારું}}.) And what do you say when you did not follow? (\textbf{\gu{માફ} \gu{કરો}}, with the words for not understanding.)
\end{tcolorbox}
